% ============================================================
% EECE 693 – Neural Networks / Deep Learning – Progress Report
% ============================================================
\documentclass[11pt,a4paper]{article}

% ── packages ──────────────────────────────────────────────────
\usepackage[margin=1in]{geometry}
\usepackage{graphicx}
\usepackage{booktabs}
\usepackage{amsmath,amssymb}
\usepackage{hyperref}
\usepackage{enumitem}
\usepackage{multirow}
\usepackage{float}
\usepackage[numbers]{natbib}
\usepackage{caption}
\usepackage{subcaption}
\usepackage{xcolor}
\usepackage{array}

\hypersetup{
  colorlinks=true,
  linkcolor=blue!60!black,
  citecolor=blue!60!black,
  urlcolor=blue!60!black
}

% ── title ─────────────────────────────────────────────────────
\title{%
  \textbf{EECE 693 (Neural Networks / Deep Learning) – Progress Report}\\[6pt]
  \Large Early Warning Prediction of Asthma Exacerbations\\
  Using Wearable and Digital-Health Data}

\author{%
  Bahaa Hamdan \qquad Tarek El Mourad \qquad Anthony Kanaan \\[4pt]
  Department of Electrical and Computer Engineering,\\
  American University of Beirut (AUB) \qquad Spring 2026\\
  Beirut, Lebanon}

\date{}

\begin{document}
\maketitle

% ============================================================
\begin{abstract}
% ============================================================
Asthma exacerbations are acute episodes of worsening respiratory symptoms that
impose a significant clinical and economic burden worldwide.  Early detection
from continuous wearable signals could enable timely intervention and reduce
hospitalizations.  In this report we describe our end-to-end pipeline
for predicting asthma exacerbations from the AAMOS-00 study (22 participants,
${\sim}2{,}054$ patient-days).  We detail (i)~data cleaning and clinically
motivated exacerbation labeling (30.3\% positive rate), (ii)~Tier-1 tabular
baselines (Logistic Regression, Random Forest, XGBoost), (iii)~Tier-2
deep-learning sequence models---a two-layer GRU (ROC-AUC\,=\,0.717) and a
CNN---and (iv)~a Tier-3 multimodal intermediate-fusion GRU that incorporates
smart-inhaler, peak-flow, and questionnaire features alongside the smartwatch
sequence.  A five-way ablation study reveals that clinical
co-modalities---especially reliever inhaler use and daily symptom
questionnaires---drive large discrimination gains: the full multimodal model
achieves ROC-AUC\,=\,0.930, F1\,=\,0.871, and Brier\,=\,0.107---a gain of
$+$0.349 over the smartwatch-only configuration in the same ablation
experiment (the standalone Tier-2 GRU achieved ROC-AUC\,=\,0.717 under full
training settings without Platt scaling).  Platt scaling on the validation
set is applied to Tier-3 multimodal test evaluations only, to mitigate the
22.5\%$\to$48.2\% train-to-test prevalence shift.
We discuss limitations including cohort size, high missing rates, and
inter-patient variability, and outline next steps including Transformer
architectures and personalisation.
\end{abstract}

\tableofcontents
\newpage

% ============================================================
\section{Introduction}\label{sec:intro}
% ============================================================

Asthma is a chronic inflammatory airway disease affecting an estimated
262~million people worldwide and contributing to approximately 455{,}000 deaths
annually~\cite{who2024}.  Exacerbations—acute or sub-acute episodes of
worsening symptoms and/or lung-function decline—represent the most
consequential events in asthma care.  They can lead to emergency visits,
hospitalization, and life-threatening deterioration if not recognized
early~\cite{bourdin2019}.  The total annual burden in the United States alone
has been estimated at~\$81.9~billion~\cite{nurmagambetov2018}.

Evidence shows that outcomes improve when worsening is identified early and
patients act on structured self-management plans~\cite{pinnock2015}.  A 2024
systematic review further supports the feasibility of detecting
exacerbation-related changes via continuous ``digital markers'' derived from
wearable and mobile-health sensors~\cite{cokorudy2024}.  Smartwatches, which
capture minute-level heart rate, activity type, step count, and intensity, are a
practical platform for such monitoring in everyday
environments~\cite{taylor2020}.

We formulate early warning as a \emph{supervised time-series prediction
problem}: given a~24-hour historical window of smartwatch signals, estimate the
probability that an exacerbation occurs within the corresponding prediction
horizon.  Using the publicly available AAMOS-00 dataset~\cite{tsang2023}, we
build an end-to-end pipeline spanning data cleaning, label construction,
feature engineering, tabular baselines, and deep-learning sequence models.  Our
contributions so far include:

\begin{enumerate}[nosep]
  \item A clinically motivated exacerbation labeling procedure combining
        documented clinical events with multi-domain symptom assessment
        (Section~\ref{sec:dataset}).
  \item Tier-1 tabular baselines (Logistic Regression, Random Forest, XGBoost)
        evaluated with patient-wise splits and rare-event metrics
        (Section~\ref{sec:results}).
  \item Tier-2 sequence models—a two-layer GRU and a three-layer 1D~CNN—trained
        on raw minute-level data (Section~\ref{sec:deep}).
  \item Threshold-tuning analysis and a GRU\,+\,RF ensemble
        (Section~\ref{sec:analysis}).
\end{enumerate}

% ============================================================
\section{Literature Review}\label{sec:litreview}
% ============================================================

\subsection{ML Landscape for Asthma Exacerbation Prediction}

A narrative review by Molfino and Turcatel~\cite{molfino2024} surveys the
machine-learning approaches applied to asthma exacerbation prediction,
highlighting that the field spans logistic regression, tree-based methods, and
deep learning.  The review emphasizes practical challenges including
heterogeneous outcome definitions, limited external validation, and the need
for generalizability across diverse patient populations.

\subsection{Home-Monitoring Time Series and the False-Alarm Problem}

De~Hond et~al.~\cite{dehond2022} compared classical and ML models for
predicting severe exacerbations from home-monitoring time series and found that
ML did not outperform logistic regression in their setting.  Importantly, they
discuss the \emph{false-alarm problem} inherent to rare-event prediction: even
models with high AUC may produce unacceptable false-positive rates when the base
rate of exacerbations is low.  Their findings motivate reporting
precision-recall AUC and sensitivity at controlled false-alarm rates, which we
adopt throughout this work.

\subsection{Connected Inhaler and Multimodal Monitoring}

Lugogo et~al.~\cite{lugogo2022} demonstrated predictive performance using
connected inhaler signals and gradient boosting, achieving clinically meaningful
discrimination.  Tsang et~al.~\cite{tsang2023} introduced the AAMOS-00
dataset, combining smartwatch sensing, smart inhaler data, peak-flow
measurements, and questionnaires, explicitly designed for asthma attack
prediction research.  The DIGIPREDICT protocol~\cite{chan2024} further reflects
the field's movement toward multimodal, physiological early-warning systems.

\subsection{Digital Markers and Wearable Data Quality}

A 2024 systematic review by Cokorudy et~al.~\cite{cokorudy2024} identifies
multiple digital markers associated with asthma exacerbations and evaluates
their potential for near real-time monitoring.  Cho et~al.~\cite{cho2021}
highlight wearable data quality challenges including missing segments and
sensor noise, motivating robust preprocessing and missingness-aware modeling.

\subsection{Summary of Gaps}

Key gaps remain in: (i)~robust modeling under rare-event conditions with
appropriate calibration and false-alarm control; (ii)~effective learning from
noisy, incomplete, heterogeneous wearable time series; and (iii)~principled
fusion of multimodal streams.  Our work addresses gaps~(i) and~(ii) in this
report and plans multimodal integration for the next phase.

% ============================================================
\section{Dataset}\label{sec:dataset}
% ============================================================

\subsection{AAMOS-00 Overview}

We use the AAMOS-00 dataset~\cite{tsang2023}, publicly available under
CC~BY~4.0~\cite{datashare}.  Phase~2 of the study enrolled 22~participants;
after cleaning and windowing, 18~users contribute usable labeled windows,
providing ${\sim}2{,}054$~unique patient-days.  In the current baseline
phase, we focus exclusively on the \textbf{smartwatch stream}, which provides
minute-by-minute recordings of four signals: heart rate (HR), activity type
(categorical), activity intensity (0--255), and step count.

\subsection{Data Cleaning}

Raw smartwatch data were loaded from three CSV files totaling all available
recording sessions.  Cleaning steps included:

\begin{itemize}[nosep]
  \item Removal of duplicate timestamps per user.
  \item Physiologically implausible HR values (${<}30$ or ${>}220$~bpm) set to
        missing.
  \item Forward-fill imputation for short gaps (${\leq}5$~min); longer gaps
        left as NaN.
  \item Users with fewer than~7 days of data retained but flagged.
\end{itemize}

After cleaning, the dataset contains \textbf{2{,}101{,}759 minute-level rows}
across \textbf{20~unique users}.  The overall HR missingness rate is
\textbf{25.3\%}, driven primarily by periods when the watch was not worn or
charging.  Step data has zero missingness.  Table~\ref{tab:usersummary}
summarizes per-user data availability.

\begin{table}[H]
\centering
\caption{Per-user data summary after cleaning (selected users shown).}
\label{tab:usersummary}
\small
\begin{tabular}{lrrr}
\toprule
\textbf{User} & \textbf{Rows} & \textbf{Days} & \textbf{HR Miss (\%)} \\
\midrule
113 & 240{,}599 & 183 & 19.2 \\
190 & 181{,}965 & 147 & 18.1 \\
294 & 224{,}323 & 184 & 19.0 \\
343 & 198{,}334 & 183 & 18.6 \\
702 & 213{,}400 & 182 & 19.3 \\
939 & 193{,}466 & 177 & 18.5 \\
808 &  83{,}193 & 154 & 61.6 \\
\bottomrule
\end{tabular}
\end{table}

\subsection{Exploratory Data Analysis}

\paragraph{Signal distributions.}
Heart rate across all users shows a mean of $81.6 \pm 9.8$~bpm with a median of
80.6~bpm.  Nighttime HR is lower (mean~$67.2$~bpm) consistent with circadian
physiology.  Activity fractions reveal that users spend on average 42.7\% of the
day sedentary, 33.0\% sleeping, 11.4\% walking, and 12.7\% in other active
states.  Mean daily step count is $6{,}604 \pm 5{,}761$ steps with high
inter-user variability.

\paragraph{Missingness patterns.}
HR missingness varies substantially across users (range: 15.8\%--61.6\%),
reflecting real-world adherence heterogeneity.  Users with $>$40\% HR
missingness (users~473, 764, 808, 917) were retained to preserve sample size
but constitute a known source of noise.  The ``charging'' and ``not worn''
activity states account for the majority of missing HR segments.

\paragraph{Temporal coverage.}
Study duration ranges from 1~day (user~748) to 184~days (user~294), with a
median of ${\sim}110$~days.  Four users contribute fewer than 15~days of data,
limiting their contribution to the model but still providing some within-user
variability.

\subsection{Labeling}

AAMOS-00 provides weekly questionnaire scores across six symptom domains
(daytime symptoms, nighttime symptoms, activity limitation, breathlessness,
wheeze, and rescue-inhaler use) together with clinical event flags (doctor
visit, hospital/ER admission, oral steroid course).  We define a week as
an exacerbation if a clinical event occurred \emph{or} at least three of
the six symptom domains are simultaneously elevated ($\geq 4$ on the 0--6
scale).  Requiring either an objectively documented clinical escalation or
concurrent multi-domain symptom worsening ensures that the positive label
captures clinically meaningful episodes rather than isolated mild complaints.
Under this rule the dataset contains \textbf{31{,}125} labeled 24-hour
windows across \textbf{18~users}, with a positive rate of
\textbf{30.3\%} (9{,}424 positive vs.\ 21{,}701 negative).  Eleven of
18~users have at least one positive window, ensuring representation across
the train/test split.

% ============================================================
\section{Windowing and Feature Engineering}\label{sec:features}
% ============================================================

\subsection{Window Construction}

We convert longitudinal data into supervised samples using:
\begin{itemize}[nosep]
  \item \textbf{History window:} 24~hours (1{,}440 minutes) of smartwatch data.
  \item \textbf{Stride:} 1~hour (60~minutes), producing overlapping windows.
  \item \textbf{Minimum coverage:} $\geq$70\% non-missing minutes required for
        a window to be retained.
\end{itemize}

Each window inherits the exacerbation label of its corresponding week.
This yields \textbf{31{,}125 windows} after coverage filtering.

\subsection{Engineered Features (Tier~1)}

For tabular baselines, we extract 29~features per window spanning four
categories:

\begin{enumerate}[nosep]
  \item \textbf{Heart rate statistics:} mean, median, standard deviation, min,
        max, slope (linear trend), and missingness percentage.
  \item \textbf{Step features:} sum, mean, max, and slope.
  \item \textbf{Activity fractions:} proportion of window spent in each activity
        state (sedentary, sleep, walk, run, high activity, charging, unknown,
        not worn).
  \item \textbf{Derived features:} active-minute fraction, day/night HR means,
        day/night step sums, day/night active fractions, intensity statistics
        (mean, std, max), and window coverage ratio.
\end{enumerate}

Feature distributions were inspected for normality and outliers.  HR standard
deviation (mean~$15.2 \pm 5.2$~bpm) and intensity standard deviation showed
the highest discriminative power in the Random Forest model
(Section~\ref{sec:results}).

% ============================================================
\section{Models}\label{sec:models}
% ============================================================

\subsection{Evaluation Protocol}

All models are evaluated using a single \textbf{patient-wise split} via
\texttt{GroupShuffleSplit} ($n_{\text{splits}}=1$, $\text{test\_size}=0.25$,
$\text{random\_state}=42$), ensuring that no patient appears in both train and
test sets.  This prevents data leakage from intra-patient correlations.  We
report seven metrics: accuracy, precision, recall, F1-score, ROC-AUC, PR-AUC,
and Brier score.

Additionally, for the Random Forest, we performed \textbf{5-fold GroupKFold
cross-validation} to estimate variance: mean AUC = $0.611 \pm 0.226$ (range:
0.37--0.91), confirming substantial inter-patient variability.

\paragraph{Class imbalance handling.}
Due to the train/test positive-rate mismatch (train: 22.5\%, test: 48.2\%—a
consequence of the small number of users combined with patient-level splitting),
we apply \texttt{class\_weight='balanced'} for Logistic Regression and
\texttt{scale\_pos\_weight} for XGBoost.  For deep-learning models, we use
\texttt{BCEWithLogitsLoss} with $\text{pos\_weight} = 5.08$ to compensate for
the training-set imbalance.

% ── 5.1 Tier 1 ──
\subsection{Tier~1: Tabular Baselines}\label{sec:tier1}

\paragraph{Logistic Regression (L2).}
Regularized logistic regression with \texttt{class\_weight='balanced'} serves
as the interpretable baseline.  It achieves moderate discrimination
(F1\,=\,0.554, ROC-AUC\,=\,0.578) but with balanced precision and recall
(52.3\%\,/\,58.8\%), indicating limited ability to separate the classes.

\paragraph{Random Forest.}
A 200-tree Random Forest with \texttt{class\_weight='balanced'} achieves the
best tabular F1 (0.547) and ROC-AUC (0.633).  It exhibits high precision
(85.4\%) but low recall (40.2\%), suggesting that it confidently identifies a
subset of exacerbation patterns but misses many positive windows.

Feature importance analysis (mean decrease in impurity) reveals that
\textbf{HR~std} (8.4\%), \textbf{night~HR~mean} (6.0\%), \textbf{night~step
sum} (6.0\%), and \textbf{intensity~std} (5.5\%) are the top predictors,
consistent with physiological expectations: elevated heart-rate variability and
nighttime disturbances may precede exacerbations.

\paragraph{XGBoost.}
XGBoost with tuned \texttt{scale\_pos\_weight} yields the highest precision
among tabular models (81.2\%) but the lowest recall (25.5\%), resulting in
F1\,=\,0.388.  Its conservative predictions limit sensitivity for early
warning.

% ── 5.2 Tier 2 ──
\subsection{Tier~2: Deep Learning on Sequences}\label{sec:deep}

Rather than using hand-crafted features, the sequence models operate on raw
minute-level data.  Each 24-hour window is represented as a
$1{,}440 \times 4$ matrix (HR, activity type encoded as integer, intensity,
steps), with NaN values zero-filled and per-feature standardization applied.

\subsubsection{GRU (2-Layer)}\label{sec:gru}

\paragraph{Architecture.}
The GRU classifier consists of a two-layer GRU with hidden dimension~64 and
dropout~0.3 between layers.  To reduce computational cost, the input sequence
is subsampled by a factor of~5 ($1{,}440 \to 288$ timesteps) before being fed
to the GRU.  The final hidden state is passed through a linear layer producing
a single logit.  Total parameters: \textbf{38{,}465}.

\paragraph{Training.}
The model is trained with Adam (lr $= 10^{-3}$), \texttt{BCEWithLogitsLoss}
(pos\_weight~$= 5.08$), batch size~64, and early stopping on validation loss
with patience~5.  Training was performed on CPU and converged in
${\sim}11$~epochs (best model at epoch~4).

\paragraph{Results.}
The GRU achieves the best overall discrimination: \textbf{ROC-AUC\,=\,0.717},
\textbf{PR-AUC\,=\,0.773}, and \textbf{F1\,=\,0.632}.  Its recall of
\textbf{91.8\%} is the highest among all models, meaning it captures the vast
majority of true exacerbation windows—a desirable property for an early-warning
system.  However, precision is moderate (48.1\%), reflecting a higher
false-alarm rate that may be tuned via threshold adjustment
(Section~\ref{sec:threshold}).

\subsubsection{1D CNN (3-Layer)}\label{sec:cnn}

\paragraph{Architecture.}
To test whether local temporal patterns suffice for exacerbation prediction, we
implement a 1D~CNN with three convolutional blocks.  Each block applies
\texttt{Conv1d} $\to$ \texttt{BatchNorm1d} $\to$ \texttt{ReLU} $\to$
\texttt{MaxPool1d(2)}, with filter counts [32, 64, 128] and kernel size~7.
Global average pooling aggregates the temporal dimension, followed by a linear
classifier with dropout~0.3.  Total parameters: \textbf{73{,}377}.  The CNN
processes the full $1{,}440$-timestep sequence (no subsampling).

\paragraph{Training.}
Same optimizer and loss as the GRU.  Training converged in ${\sim}10$~epochs,
with best validation loss significantly higher than the GRU.

\paragraph{Results.}
The CNN achieves near-random discrimination: \textbf{ROC-AUC\,=\,0.520},
\textbf{PR-AUC\,=\,0.530}, \textbf{F1\,=\,0.564}.  This result is
substantially worse than the GRU across all metrics.  The CNN primarily
captures short-term temporal patterns such as transient heart-rate spikes
and brief activity bursts, whereas exacerbation onset is a gradual process
spanning hours to days.  Given the small dataset size, we intentionally used
a lightweight three-block architecture to limit overfitting; deeper or
dilated designs (e.g.\ TCN) could extend the receptive field but risk
increased variance.  The GRU's recurrent architecture, even with subsampled
input, maintains state across the full 24-hour context and models the
long-range dependencies that prove critical for this task.

% ============================================================
\section{Results}\label{sec:results}
% ============================================================

\subsection{Full Model Comparison}

Table~\ref{tab:results} presents all seven models evaluated on the held-out
patient-wise test set, sorted by ROC-AUC.

\begin{table}[H]
\centering
\caption{Performance comparison of all models on the patient-wise test set
(sorted by ROC-AUC). Best values per metric in \textbf{bold}.}
\label{tab:results}
\small
\begin{tabular}{lccccccc}
\toprule
\textbf{Model} & \textbf{Acc} & \textbf{Prec} & \textbf{Rec} & \textbf{F1}
  & \textbf{ROC-AUC} & \textbf{PR-AUC} & \textbf{Brier} \\
\midrule
GRU (2-layer)       & 0.483 & 0.481 & \textbf{0.918} & 0.632 & \textbf{0.717} & \textbf{0.773} & 0.225 \\
Ensemble (60\%\,GRU) & 0.726 & 0.823 & 0.549 & \textbf{0.659} & 0.653 & 0.762 & \textbf{0.221} \\
Ensemble (avg)      & \textbf{0.726} & \textbf{0.839} & 0.535 & 0.654 & 0.651 & 0.758 & 0.223 \\
Random Forest       & 0.678 & 0.854 & 0.402 & 0.547 & 0.633 & 0.724 & 0.254 \\
XGBoost             & 0.612 & 0.812 & 0.255 & 0.388 & 0.581 & 0.654 & 0.312 \\
Logistic Reg.       & 0.543 & 0.523 & 0.588 & 0.554 & 0.578 & 0.658 & 0.313 \\
CNN (3-layer)       & 0.483 & 0.475 & 0.694 & 0.564 & 0.520 & 0.530 & 0.320 \\
\bottomrule
\end{tabular}
\end{table}

\subsection{Key Observations}

\begin{enumerate}[nosep]
  \item \textbf{GRU dominates in discrimination.}  The two-layer GRU achieves
        the highest ROC-AUC (0.717) and PR-AUC (0.773), outperforming the best
        tabular model (RF ROC-AUC\,=\,0.633) by 8.4 percentage points.  For
        reference, the no-skill PR-AUC baseline equals the positive-class
        prevalence in the test set (0.482), so all models exceed chance.  This
        confirms that end-to-end learning from raw sequences captures
        predictive patterns missed by hand-crafted features.

  \item \textbf{Ensemble improves calibration and precision.}  Combining GRU
        probabilities with RF predictions (60\%\,GRU\,+\,40\%\,RF) yields the
        best F1 (0.659) and Brier score (0.221), improving precision from
        48.1\% to 82.3\% at the cost of recall (91.8\%\,$\to$\,54.9\%).

  \item \textbf{CNN fails on sequence-only data.}  Near-random AUC
        (0.520) indicates that local convolutional patterns do not capture the
        relevant temporal dynamics under full training settings.  The CNN's
        receptive field, even with three pooling layers, may be too narrow to
        model the slow-evolving physiological changes preceding exacerbations.

  \item \textbf{Sequence-only models fall far below multimodal fusion.}
        The smartwatch-only GRU (AUC\,=\,0.581 in the ablation experiment)
        scores $-$0.349 below the full multimodal GRU (AUC\,=\,0.930),
        confirming that clinical co-modalities—not architectural
        complexity—are the primary driver of discrimination gain.

  \item \textbf{Tabular baselines are moderate but limited.}  All three
        Tier-1 models cluster between AUC~0.578 and 0.633, with engineered
        features providing a useful but incomplete representation.

  \item \textbf{High inter-patient variance.}  5-fold GroupKFold CV for the RF
        shows AUC ranging from 0.37 to 0.91, confirming that exacerbation
        patterns vary substantially between individuals.  This motivates
        personalized or multi-task modeling in future work.
\end{enumerate}

% ============================================================
\section{Model Analysis and Robustness}\label{sec:analysis}
% ============================================================

\subsection{Threshold Tuning}\label{sec:threshold}

Given class imbalance and asymmetric clinical costs—missed exacerbations
carry greater risk than false alarms—we evaluate model performance across
decision thresholds rather than relying solely on a fixed threshold of~0.5.
The GRU's default threshold maximizes recall (91.8\%) but at the cost of
precision (48.1\%).  We perform systematic threshold analysis to identify
operating points suitable for different clinical scenarios:

\begin{table}[H]
\centering
\caption{GRU operating points at different decision thresholds.}
\label{tab:thresholds}
\small
\begin{tabular}{lcccc}
\toprule
\textbf{Operating Point} & \textbf{Threshold} & \textbf{Precision} &
  \textbf{Recall} & \textbf{F1} \\
\midrule
Default (0.50)            & 0.500 & 0.481 & 0.918 & 0.632 \\
Best F1                   & 0.525 & 0.539 & 0.798 & 0.653 \\
FPR-controlled            & 0.531 & 0.719 & 0.550 & 0.623 \\
\bottomrule
\end{tabular}
\end{table}

The Best-F1 threshold (0.525) achieves a slight improvement in F1 (0.653 vs.\
0.632) by trading recall (91.8\% $\to$ 79.8\%) for precision (48.1\% $\to$
53.9\%).  The FPR-controlled threshold (0.531) dramatically improves precision
to 71.9\% while maintaining clinically useful recall (55.0\%).  In a
real-world deployment, the choice of threshold would depend on whether the
system prioritizes sensitivity (to catch all exacerbations) or specificity (to
minimize alarm fatigue).

\subsection{Feature Importance Analysis}

The Random Forest's top-5 features by mean decrease in impurity are:

\begin{enumerate}[nosep]
  \item HR standard deviation (8.4\%) — captures heart-rate variability
        changes.
  \item Night HR mean (6.0\%) — elevated nighttime HR may signal respiratory
        distress.
  \item Night step sum (6.0\%) — increased nighttime movement suggesting
        sleep disruption.
  \item Intensity standard deviation (5.5\%) — variability in activity
        intensity.
  \item Intensity max (5.1\%) — peak activity-intensity level.
\end{enumerate}

The dominance of HR variability and nighttime features is consistent with
clinical literature: nocturnal symptoms and autonomic changes are established
markers of worsening asthma control~\cite{bourdin2019}.

\subsection{Ensemble Analysis}

We evaluated two ensemble strategies combining GRU and Random Forest
predictions on the same test set:

\begin{itemize}[nosep]
  \item \textbf{Average ensemble} (50\%\,GRU + 50\%\,RF): F1\,=\,0.654,
        AUC\,=\,0.651, Brier\,=\,0.223.
  \item \textbf{Weighted ensemble} (60\%\,GRU + 40\%\,RF): F1\,=\,0.659,
        AUC\,=\,0.653, Brier\,=\,\textbf{0.221} (best calibration).
\end{itemize}

The ensemble benefits from model diversity: the GRU captures temporal
dynamics from raw minute-level sequences, while the RF operates on aggregated
statistical features, producing complementary error patterns.  Probability
averaging smooths extreme predictions from either model and reduces
overconfidence, which is reflected in the improved Brier scores.
Both ensembles trade GRU recall for RF precision, producing models with
balanced precision–recall profiles.  The weighted ensemble achieves the best
Brier score across all models, indicating superior probability calibration.
The 60/40 weighting reflects the GRU's superior standalone AUC rather than
being optimized on the test set.  However, the ensembles do not surpass the
standalone GRU in terms of discrimination (AUC), consistent with the
observation that combining a strong model with a weaker one can dilute ranking
performance while improving point-estimate quality.

\subsection{Robustness Concerns}

\paragraph{Small cohort and single split.}
With only 18~patients (the split placing ${\sim}4$--5 in the test set),
results are sensitive to which patients happen to be in train vs.\ test.
All reported metrics come from a single patient-wise split (seed\,=\,42).
Bootstrap 95\% confidence intervals ($n{=}1{,}000$ resamples, seed\,=\,42)
are computed via percentile bootstrap on the test set.  A locally reproducible
run (GRU seed\,=\,8) yields GRU AUC\,=\,0.665 [0.655,\,0.675] and
CNN AUC\,=\,0.582 [0.572,\,0.593].  The within-run CI width is
${\sim}0.020$, but inter-run variation due to initialisation sensitivity
(GRU spans 0.394--0.665 across random seeds) exceeds this, so differences
smaller than ${\sim}0.05$ are not statistically distinguishable.  The
originally reported GRU AUC of 0.717 was obtained on Colab and reflects a
favourable initialisation; the ${\sim}0.05$ gap relative to the local
result is within this noise floor.

\paragraph{Effective sample size.}
Adjacent windows share 95.8\% of their data (24-hour window, 1-hour stride).
The reported ${\sim}31{,}000$ windows correspond to approximately
${\sim}1{,}900$ \emph{independent} observations.  All reported model
comparisons and significance claims should be interpreted in light of
this effective sample size, not the nominal window count.

\paragraph{Train/test prevalence shift.}
The patient-wise split results in different positive rates between train
(22.5\%) and test (48.2\%) sets---a consequence of the small cohort and
patient-level splitting.  This shift inflates PR-AUC, Brier score, and
F1 on the test set relative to a balanced deployment setting.  ROC-AUC
is approximately invariant to prevalence and is the recommended primary
metric.  Platt scaling is applied only to Tier-3 multimodal models; Tier-1
and Tier-2 evaluations use raw model probabilities, introducing a calibration
inconsistency across tiers.

\paragraph{Cross-validation on the full dataset.}
The 5-fold GroupKFold cross-validation for Tier-1 models was run on the
full dataset including test-set patients.  The resulting fold AUC estimates
(range: 0.37--0.91) reflect inter-patient variance only and are not valid
held-out performance estimates; they should not be compared directly with the
held-out test results.

\paragraph{Cross-protocol comparisons.}
The Tier-2 GRU (full-run: subsample\,=\,5, 20~epochs, no Platt scaling,
ROC-AUC\,=\,0.717) and the Tier-3 multimodal ablation (fast-run:
subsample\,=\,10, 10~epochs, with Platt scaling) were trained under
different settings.  These protocols are not directly comparable; the
$+$0.349 gain reported in the ablation is a within-experiment comparison
(\texttt{Full} vs.\ \texttt{SW\_only}).

\paragraph{Multimodal reproducibility.}
The raw co-modality data files (smart inhaler logs, peak-flow records,
daily questionnaire responses, weekly questionnaire scores) are not
included in the local repository.  Tier-3 multimodal results cannot be
reproduced without access to the full AAMOS-00 dataset release from the
original data repository.

\paragraph{Window overlap.}
Adjacent windows (1-hour stride over 24-hour duration) share 95.8\% of their
data, creating redundancy within each user.  This does not leak across the
train/test boundary (patient-wise split), but inflates effective sample counts
and may bias variance estimates.

% ============================================================
\section{Efficiency and Scalability}\label{sec:efficiency}
% ============================================================

\begin{table}[H]
\centering
\caption{Model complexity and training efficiency (CPU only).}
\label{tab:efficiency}
\small
\begin{tabular}{lrrr}
\toprule
\textbf{Model} & \textbf{Parameters} & \textbf{Input Shape} &
  \textbf{Train Time} \\
\midrule
Logistic Regression & 30 &  $29$ features & $<$1\,s \\
Random Forest       & -- (200 trees) & $29$ features & ${\sim}5$\,s \\
XGBoost             & -- (100 trees) & $29$ features & ${\sim}3$\,s \\
GRU (2-layer)       & 38{,}465 & $288 \times 4$ & ${\sim}5$\,min \\
CNN (3-layer)       & 73{,}377 & $1{,}440 \times 4$ & ${\sim}4.8$\,min \\
\bottomrule
\end{tabular}
\end{table}

All experiments were conducted on CPU (no GPU acceleration).  The GRU's
subsampling strategy (factor of~5) reduces the effective sequence length from
1{,}440 to 288 without degrading performance—in fact, the GRU outperforms the
CNN which processes the full sequence.  This suggests that 5-minute resolution
captures sufficient temporal dynamics for exacerbation prediction, an important
finding for potential deployment on resource-constrained wearable devices.

The tabular baselines train in seconds but require a separate feature
engineering step; the deep models train in minutes but operate end-to-end.  For
real-time inference, both the GRU (38K parameters) and CNN (73K parameters) are
small enough for edge deployment on modern smartwatch processors.

% ============================================================
\section{Multimodal Fusion Modeling}\label{sec:multimodal}
% ============================================================

\subsection{Architecture}

We implement an \textbf{intermediate-fusion GRU} that combines the
smartwatch minute-level sequence with tabular features derived from clinical
co-modalities (smart inhaler, peak-flow meter, and questionnaires).

\begin{figure}[H]
\centering
\small
\begin{tabular}{l}
\textbf{Sequence branch} \\
Smartwatch: (batch, 288, 4) $\xrightarrow{\text{GRU}(h{=}64,\,L{=}2)}$ $h_T$ \in $\mathbb{R}^{64}$ \\[4pt]
\textbf{Tabular branch} (omitted when $K=0$) \\
Clinical features: (batch, $K$) $\xrightarrow{K \to 64 \to 32,\;\text{ReLU, Dropout}}$ $\mathbf{z}_\text{tab}$ \in $\mathbb{R}^{32}$ \\[4pt]
\textbf{Fusion} \\
$[\,h_T\,\|\,\mathbf{z}_\text{tab}\,] \in \mathbb{R}^{96}$ (or $\mathbb{R}^{64}$ for SW-only) $\xrightarrow{\text{Linear}(1)}$ logit
\end{tabular}
\caption{Intermediate-fusion GRU architecture. The sequence is subsampled by
$5\times$ (1\,440 $\to$ 288 steps) to reduce memory and compute.}
\label{fig:fusion_arch}
\end{figure}

\noindent\textbf{Training details.}
BCEWithLogitsLoss with $\text{pos\_weight}=3.43$ (computed from training-set
class counts) is used for all configs to address the 22.5\% training positive
rate.  Adam optimiser ($\text{lr}=10^{-3}$, $\text{weight\_decay}=10^{-4}$)
with early stopping (patience\,=\,5) on validation loss.  Platt scaling is
fitted on validation-set probabilities and applied to all test-set
evaluations to correct the train$\to$test prevalence shift
(22.5\%$\to$48.2\%).

\subsection{Tabular Feature Construction}

$K=90$ tabular features are built from three co-modalities:

\begin{itemize}[nosep]
  \item \textbf{Smart inhaler} (6 features): daily reliever-use count,
        nocturnal-use flag, inhaler coverage days, and binary missing
        indicators.
  \item \textbf{Peak-flow meter} (15 features): AM/PM PEF, morning dip,
        week-level PEF statistics, and missing indicators.
  \item \textbf{Daily questionnaire} (3 features): composite worsening score
        and temporal recency encoding.
  \item \textbf{Weekly questionnaire + patient static} (66 features):
        ACQ/AQLQ domain scores, demographic indicators, and binary flags
        for all missing entries.
\end{itemize}

All tabular features are z-score-standardised using training-set statistics
only.  Missing values are filled with sentinel $-1$ prior to standardisation;
a paired binary indicator column flags each imputed entry.

\subsection{Ablation Experiment Design}

Five cumulative ablation configurations are evaluated, each adding one
additional co-modality group to the smartwatch GRU:

\begin{table}[H]
\centering
\caption{Ablation configurations and tabular feature dimensions.}
\label{tab:ablation_configs}
\small
\begin{tabular}{llr}
\toprule
\textbf{Config} & \textbf{Input modalities} & $K$ \\
\midrule
\texttt{SW\_only}               & Smartwatch sequence only           &  0 \\
\texttt{SW\_inhaler}            & + Smart inhaler daily features     &  6 \\
\texttt{SW\_inhaler\_pef}       & + Peak-flow meter features         & 21 \\
\texttt{SW\_inhaler\_pef\_dailyQ} & + Daily questionnaire features   & 24 \\
\texttt{Full}                   & + Weekly questionnaire + static    & 90 \\
\bottomrule
\end{tabular}
\end{table}

\subsection{Results}

Table~\ref{tab:multimodal_results} reports test-set performance for all five
configurations.  \textbf{ROC-AUC is the primary metric} because it is
threshold-independent and robust to the train$\to$test prevalence shift.
PR-AUC and Brier score are reported for completeness but are inflated by the
48.2\% test-set positive rate (vs.\ 22.5\% in training).

\begin{table}[H]
\centering
\caption{Multimodal ablation results on the held-out test set (13\,539
windows, 18 patients, Platt-calibrated). F1 computed at $\tau{=}0.5$; for
\texttt{SW\_only}–\texttt{SW\_inhaler\_pef} the model predicts all-negative at
this threshold due to the prevalence shift (see text).}
\label{tab:multimodal_results}
\small
\begin{tabular}{lrrrrrrr}
\toprule
\textbf{Config} & $K$ &
  \textbf{ROC-AUC} & \textbf{PR-AUC} & \textbf{F1} &
  \textbf{Precision} & \textbf{Recall} & \textbf{Brier} \\
\midrule
\texttt{SW\_only}               & 0  & 0.581 & 0.552 & 0.000 & --- & 0.000 & 0.426 \\
\texttt{SW\_inhaler}            & 6  & 0.791 & 0.831 & 0.000 & --- & 0.000 & 0.399 \\
\texttt{SW\_inhaler\_pef}       & 21 & 0.750 & 0.801 & 0.000 & --- & 0.000 & 0.399 \\
\texttt{SW\_inhaler\_pef\_dailyQ} & 24 & 0.913 & 0.933 & 0.688 & 0.989 & 0.528 & 0.205 \\
\texttt{Full}                   & 90 & \textbf{0.930} & \textbf{0.954} & \textbf{0.871} & 0.997 & 0.774 & \textbf{0.107} \\
\midrule
GRU baseline (Tier-2)           &  0 & 0.717 & 0.773 & 0.645 & 0.673 & 0.619 & 0.220 \\
\bottomrule
\end{tabular}
\end{table}

\noindent\textbf{Key findings.}

\begin{itemize}[nosep]
  \item \textbf{Smartwatch alone is insufficient.}  The sequence-only GRU
        (\texttt{SW\_only}) achieves ROC-AUC\,=\,0.581, only marginally above
        chance, confirming that physiological signals from wearables alone are
        insufficient for reliable exacerbation prediction in this cohort.

  \item \textbf{Smart-inhaler use is the strongest single co-modality.}
        Adding only six inhaler features jumps ROC-AUC to 0.791
        ($+$0.210), demonstrating that reliever-medication behaviour is a
        clinically actionable leading indicator of impending exacerbation.

  \item \textbf{Peak-flow features add limited marginal value alone.}
        \texttt{SW\_inhaler\_pef} (ROC-AUC\,=\,0.750) dips slightly from the
        inhaler-only config, likely due to the 42.6\% missing rate for
        peak-flow readings degrading the benefit of the additional features at
        this training budget.

  \item \textbf{Daily questionnaires provide the largest single-step gain.}
        Adding the composite symptom-worsening score (\texttt{dailyQ},
        $K{=}3$ extra features) raises ROC-AUC from 0.750 to
        0.913\,($+$0.163), despite an 84.4\% missing rate, because the
        entries that \emph{are} present are highly discriminative.

  \item \textbf{Full multimodal fusion achieves ROC-AUC\,=\,0.930.}
        The complete model ($K{=}90$) yields ROC-AUC\,=\,0.930,
        PR-AUC\,=\,0.954, F1\,=\,0.871, and Brier\,=\,0.107—a substantial
        improvement of $+$0.349 ROC-AUC over the smartwatch-only baseline
        (\texttt{SW\_only}, AUC\,=\,0.581) in the same ablation experiment.
        For reference, the separately trained Tier-2 GRU achieved
        ROC-AUC\,=\,0.717 under full training settings and no Platt scaling;
        the two settings are not directly comparable.

  \item \textbf{F1 collapse for low-modality configs.}
        \texttt{SW\_only}, \texttt{SW\_inhaler}, and \texttt{SW\_inhaler\_pef}
        predict all-negative at $\tau{=}0.5$, driven by the 22.5\%$\to$48.2\%
        prevalence shift between training and test sets.  Platt scaling
        partially corrects calibration but does not fully bridge a
        2$\times$ prevalence gap for models with weaker discriminative power.
        ROC-AUC remains the appropriate comparison metric.
\end{itemize}

\subsection{Limitations}

\begin{itemize}[nosep]
  \item \textbf{Prevalence shift.}  The patient-wise split produces a
        training positive rate of 22.5\% and test positive rate of 48.2\%.
        This imbalance inflates PR-AUC and Brier on the test set and
        suppresses F1 at $\tau{=}0.5$ for under-powered configs.

  \item \textbf{High missing rates.}  Smart inhaler (63.1\%) and daily
        questionnaire (84.4\%) data are frequently missing, limiting the
        number of windows with true clinical signal.  Sentinel imputation
        with missing indicators is employed but may introduce noise.

  \item \textbf{Small cohort.}  Only 18 patients contribute labeled windows,
        making it impossible to assess generalisation beyond this cohort.

  \item \textbf{Environment data excluded.}  Weather, pollen, and air
        quality data available in AAMOS-00 are not yet incorporated.

  \item \textbf{Reduced training budget.}  The local CPU run used
        subsample\,=\,10 (144 time-steps) and 10 epochs (vs.\ 5 and 20 in
        the Tier-2 baseline), which suppresses the \texttt{SW\_only}
        baseline; the relative gains from co-modalities remain valid.

  \item \textbf{No uncertainty quantification.}  All results derive from
        a single train/val/test split.  Bootstrap confidence intervals or
        repeated-seed evaluation would be required to establish statistical
        significance of reported differences; these are left for future work.
\end{itemize}

% ============================================================
\section{Ethical Considerations}\label{sec:ethics}
% ============================================================

\paragraph{Data privacy.}
AAMOS-00 is a fully anonymized, publicly available dataset released under
CC~BY~4.0.  All user identifiers in this work are arbitrary integer keys with
no link to real identities.  No attempt is made to re-identify participants.

\paragraph{Clinical safety.}
An early-warning system for asthma exacerbations must balance sensitivity
(catching true exacerbations) against specificity (avoiding alarm fatigue).
False negatives carry higher clinical risk—missed exacerbations can escalate
rapidly.  Our threshold-tuning analysis (Section~\ref{sec:threshold}) provides
operating points at different precision–recall trade-offs, allowing clinicians
to configure the system according to the patient's risk profile and tolerance
for alerts.

\paragraph{Bias and fairness.}
The AAMOS-00 cohort (22~participants enrolled; 18 contribute labeled windows;
from a single UK-based study) is small and
may not represent the broader asthma population in terms of demographics,
disease severity, or device adherence patterns.  Models trained on this cohort
should not be deployed without validation on more diverse populations.

\paragraph{Informed consent and intended use.}
Any real-world deployment would require regulatory approval, validated clinical
trials, and integration with clinical decision support frameworks rather than
autonomous decision-making.  Our work is intended for research purposes only.

% ============================================================
\section{Conclusion and Future Work}\label{sec:conclusion}
% ============================================================

\subsection{Summary of Progress}

We have built an end-to-end pipeline for asthma exacerbation prediction,
progressing from data cleaning through labeling, feature engineering,
tabular baselines, deep-learning sequence models, and multimodal fusion.
Key findings:

\begin{itemize}[nosep]
  \item \textbf{Clinically grounded labeling.}  Combining clinical event flags
        with multi-domain symptom thresholds yields a 30.3\% positive rate,
        providing a meaningful and realistic classification challenge.
  \item \textbf{GRU excels at sequence modeling.}  The two-layer GRU achieves
        the best Tier-2 discrimination (ROC-AUC\,=\,0.717, PR-AUC\,=\,0.773)
        by learning directly from raw minute-level sequences.
  \item \textbf{Sequence-only models are insufficient.}  The smartwatch-only
        GRU achieves AUC\,=\,0.581 in the multimodal ablation experiment,
        confirming that clinical co-modalities—not architectural
        complexity—drive discrimination gains.
  \item \textbf{Ensembles improve calibration.}  The GRU\,+\,RF weighted
        ensemble achieves the best Tier-2 Brier score (0.221) and highest F1 (0.659).
  \item \textbf{Multimodal fusion achieves major gains.}  The full
        intermediate-fusion GRU ($K{=}90$ tabular features) reaches
        ROC-AUC\,=\,0.930, F1\,=\,0.871, and Brier\,=\,0.107—a gain of
        $+$0.349 ROC-AUC over the smartwatch-only baseline in the same
        ablation experiment.
  \item \textbf{Reliever use is the strongest leading indicator.}  Adding
        six smart-inhaler features alone raises ROC-AUC from 0.581 to 0.791
        ($+$0.210), the largest single-modality gain.
  \item \textbf{Questionnaire data is highly discriminative despite sparsity.}
        Daily questionnaire features (84.4\% missing) drive ROC-AUC from
        0.750 to 0.913 ($+$0.163) when combined with inhaler and peak-flow.
  \item \textbf{Inter-patient variability is a primary challenge.}  CV AUC
        ranges from 0.37 to 0.91 across patient folds, motivating
        personalization strategies.
\end{itemize}

\subsection{Planned Next Steps}

\begin{enumerate}[nosep]
  \item \textbf{Threshold optimisation.}  Apply F1-optimal and
        clinical-operating-point thresholds for the Full multimodal model to
        improve F1 beyond the $\tau{=}0.5$ default, addressing the
        train$\to$test prevalence shift.
  \item \textbf{Uncertainty quantification.}  Compute bootstrap confidence
        intervals for key AUC results and evaluate under multiple random seeds
        to characterise result stability given the small cohort.
  \item \textbf{Tier~3: Transformer models.}  Evaluate Transformer-based
        architectures (e.g.\ PatchTST, Informer) for longer temporal context
        and multi-horizon prediction.
  \item \textbf{Multiple prediction horizons.}  Evaluate 6-hour and 12-hour
        windows alongside the current 24-hour setting.
  \item \textbf{Personalisation.}  Explore user-specific fine-tuning or
        multi-task learning to reduce inter-patient CV variance.
  \item \textbf{Environment data integration.}  Incorporate weather, pollen,
        and air quality features available in AAMOS-00.
  \item \textbf{Missing-modality robustness.}  Investigate masked-autoencoder
        or cross-attention strategies to better exploit partially observed
        co-modalities with high missing rates.
\end{enumerate}

% ============================================================
% References
% ============================================================
\bibliographystyle{unsrtnat}

\begin{thebibliography}{20}

\bibitem{who2024}
World Health Organization, ``Asthma,'' Fact Sheet, May 6, 2024.

\bibitem{bourdin2019}
A.~Bourdin et~al., ``ERS/EAACI statement on severe exacerbations in asthma,''
\emph{European Respiratory Journal}, vol.~54, no.~3, 2019.

\bibitem{nurmagambetov2018}
T.~Nurmagambetov, R.~Kuwahara, and P.~Garbe, ``The economic burden of asthma
in the United States, 2008--2013,'' \emph{Annals of the American Thoracic
Society}, 2018.

\bibitem{pinnock2015}
H.~Pinnock, ``Supported self-management for asthma,'' \emph{Breathe (ERS)},
2015.

\bibitem{cokorudy2024}
B.~Cokorudy et~al., ``Digital markers of asthma exacerbations: a systematic
review,'' 2024.

\bibitem{taylor2020}
L.~Taylor et~al., ``Wearable vital signs monitoring for patients with asthma:
a review,'' 2020.

\bibitem{cho2021}
J.~Cho et~al., ``Wearable data quality: a systematic review,'' 2021.

\bibitem{molfino2024}
N.~A.~Molfino and G.~Turcatel, ``Machine learning approaches to predict asthma
exacerbations: a narrative review,'' 2024.

\bibitem{dehond2022}
A.~A.~H.~de~Hond et~al., ``Machine learning did not beat logistic regression
in time series prediction for severe asthma exacerbations,'' \emph{Scientific
Reports}, vol.~12, Art.~no.~20363, 2022.

\bibitem{lugogo2022}
N.~L.~Lugogo et~al., ``A predictive machine learning tool for asthma
exacerbations using an electronic inhaler with integrated sensors,''
\emph{Journal of Asthma and Allergy}, 2022.

\bibitem{tsang2023}
K.~C.~H.~Tsang et~al., ``Home monitoring with connected mobile devices for
asthma attack prediction with machine learning,'' \emph{Scientific Data},
2023.

\bibitem{chan2024}
A.~H.~Y.~Chan et~al., ``DIGIPREDICT: physiological, behavioural and
environmental predictors of asthma attacks---study protocol,'' \emph{BMJ Open
Respiratory Research}, 2024.

\bibitem{datashare}
University of Edinburgh DataShare, ``AAMOS-00 dataset repository,'' CC~BY~4.0.

\bibitem{camargo2009}
C.~A.~Camargo~Jr., G.~Rachelefsky, and M.~Schatz, ``Managing asthma
exacerbations in the emergency department,'' \emph{Proc.\ Am.\ Thorac.\
Soc.}, 2009.

\bibitem{ke2017}
G.~Ke et~al., ``LightGBM: A highly efficient gradient boosting decision
tree,'' \emph{NeurIPS}, 2017.

\bibitem{bai2018}
S.~Bai, J.~Z.~Kolter, and V.~Koltun, ``An empirical evaluation of generic
convolutional and recurrent networks for sequence modeling,'' 2018.

\bibitem{zhou2021}
H.~Zhou et~al., ``Informer: Beyond efficient transformer for long sequence
time-series forecasting,'' \emph{AAAI}, 2021.

\bibitem{lim2021}
B.~Lim et~al., ``Temporal fusion transformers for interpretable multi-horizon
time series forecasting,'' \emph{International Journal of Forecasting}, 2021.

\bibitem{ganaie2021}
M.~A.~Ganaie et~al., ``Ensemble deep learning: a review,'' 2021.

\bibitem{synapse}
Synapse, ``Asthma Mobile Health Study (AMHS),'' dataset record.

\end{thebibliography}

\end{document}
